% answers_short.tex —— 2020 年·全国Ⅱ卷·文科 · 简短答案速查
% 由 tools/make_short_answers.py 自 answers.tex 的 \gksolhead/\gkanswer 提取，
% 并与 2020/2020_数学真题_全国2文_答案要点.md 逐题比对后生成，勿手改（改源 MD 后重跑）。

\begin{gknotice}
\textbf{使用说明}：本节只给最终答案便于快速核对，完整推导见解析版。答案为编写者独立推导，\textbf{非官方标准答案}。
\end{gknotice}

\gkpart{选择题}
\begin{center}\begin{tabular}{ccccccccccccc}
\toprule
题号 & 1 & 2 & 3 & 4 & 5 & 6 & 7 & 8 & 9 & 10 & 11 & 12 \\
\midrule
答案 & D & A & C & B & D & B & C & B & B & A & C & A \\
\bottomrule
\end{tabular}\end{center}

\gkpart{填空题}
\begin{itemize}[leftmargin=2.6em]
  \item[\textbf{13.}] $\dfrac19$。$\cos2x=1-2\sin^2x=1-2\cdot\dfrac49=\dfrac19$。
  \item[\textbf{14.}] $25$。由 $a_2+a_6=2a_1+6d=2$、$a_1=-2$ 得 $d=1$，故 $S_{10}=10a_1+45d=25$。
  \item[\textbf{15.}] $8$。可行域的三个顶点为 $(-1,0)$、$(0,-1)$、$(2,3)$，$z=x+2y$ 在 $(2,3)$ 处取最大值 $8$。
  \item[\textbf{16.}] ①③④。$p_1$、$p_4$ 为真，$p_2$、$p_3$ 为假，代入四种复合形式即得。
\end{itemize}

\gkpart{解答题}
\begin{itemize}[leftmargin=2.6em]
  \item[\textbf{17.}] (1) $A=\dfrac\pi3$。由 $\cos\left(\dfrac\pi2+A\right)=-\sin A$ 得 $\sin^2A+\cos A=\dfrac54$，即 $\left(\cos A-\dfrac12\right)^2=0$；(2) 证明。由余弦定理与 $a=\sqrt3(b-c)$ 得 $(2b-c)(b-2c)=0$；$c=2b$ 时 $c^2=a^2+b^2$，$b=2c$ 时 $b^2=a^2+c^2$，故必为直角三角形
  \item[\textbf{18.}] (1) $12000$。样本平均数 $\bar y=\dfrac{1200}{20}=60$，乘地块数 $200$ 得 $60\times200=12000$；(2) $0.94$。$r=\dfrac{800}{\sqrt{80\times9000}}=\dfrac{2\sqrt2}3\approx0.94$；(3) 分层抽样（按植物覆盖面积分层）。地块间差异大而层内差异小，分层后样本结构与总体一致，估计更准确
  \item[\textbf{19.}] (1) $e=\dfrac12$。由 $|CD|=\dfrac43|AB|$ 得 $4c=\dfrac43\cdot\dfrac{2b^2}a$，即 $2a^2-3ac-2c^2=0$，解 $2e^2+3e-2=0$；(2) $C_1:\dfrac{x^2}{16}+\dfrac{y^2}{12}=1$，$C_2:y^2=8x$。由 $a=2c$ 及四个顶点到准线 $x=-c$ 的距离之和 $2a+2c=12$ 得 $c=2$
  \item[\textbf{20.}] (1) 证明。$MN\parallel BB_1\parallel AA_1$；又 $BC\perp AM$、$BC\perp MN$，故 $BC\perp$ 平面 $A_1AMN$，而 $B_1C_1\parallel BC$ 且 $B_1C_1\subset$ 平面 $EB_1C_1F$，故两平面垂直；(2) $24$。由 $AO\parallel$ 平面 $EB_1C_1F$ 得 $AP=\dfrac13AM$，故 $AE=AF=2$，$EF=2$；再由 $\angle MPN=\dfrac\pi3$ 与 $AO=6$ 定出棱柱的高，得四棱锥高为 $3$、底面梯形面积为 $24$
  \item[\textbf{21.}] (1) $c\geqslant-1$。令 $h(x)=2\ln x+1-2x$，$h'(x)=\dfrac2x-2$，最大值 $h(1)=-1$；(2) $g(x)$ 在 $(0,a)$ 上单调递减，在 $(a,+\infty)$ 上单调递减。由 $g'(x)=\dfrac{2\left[1-\frac ax-\ln\frac xa\right]}{(x-a)^2}$ 及 $1-\dfrac1t-\ln t\leqslant0$ 定号
\end{itemize}

\gkpart{选做题}
\begin{itemize}[leftmargin=2.6em]
  \item[\textbf{22.}] (1) $C_1:x+y=4$（$0\leqslant x\leqslant4$），$C_2:x^2-y^2=4$。前者由 $4\cos^2\theta+4\sin^2\theta=4$ 消参，后者由 $(x+y)(x-y)=4$ 消参；(2) $\rho=\dfrac{17}5\cos\theta$。交点 $P\left(\dfrac52,\dfrac32\right)$，圆心 $\left(\dfrac{17}{10},0\right)$，故 $\rho=2\cdot\dfrac{17}{10}\cos\theta$
  \item[\textbf{23.}] (1) $\left(-\infty,\dfrac32\right]\cup\left[\dfrac{11}2,+\infty\right)$。$a=2$ 时 $f(x)=|x-4|+|x-3|$，分段求解 $f(x)\geqslant4$；(2) $a\leqslant-1$ 或 $a\geqslant3$。$f(x)$ 的最小值为 $(a-1)^2$，故 $(a-1)^2\geqslant4$
\end{itemize}
